\chapter{WebNLG: ediciones, estructura del \textit{dataset} y extensión al español} \label{cha:webnlg-desarrollo}

La organización oficial de WebNLG se encuentra disponible en el repositorio
público de GitHub\footnote{\url{https://github.com/webnlg}}, donde se alojan
los recursos empleados en las distintas ediciones del \textit{challenge}. A lo
largo de los años, WebNLG ha sido utilizado como \textit{benchmark} estándar
para la evaluación comparativa de sistemas de generación y estructuración de
lenguaje natural.

\section{Ediciones del WebNLG Challenge}

\subsection{WebNLG Challenge 2017}

El WebNLG Challenge 2017 constituyó la primera competición en la que el
conjunto de datos WebNLG fue empleado formalmente como referencia para evaluar
sistemas de generación de texto a partir de tripletas RDF. La tarea era
exclusivamente \textit{data-to-text} y el idioma considerado fue el inglés.

Las métricas de evaluación empleadas fueron BLEU, METEOR y TER. El desafío fue
presentado en la 10th International Conference on Natural Language Generation
(INLG 2017), celebrada en Santiago de Compostela~\cite{Gardent2017}.

\subsection{WebNLG+ Challenge 2020}

La edición WebNLG+ 2020 amplió el alcance del desafío incorporando una tarea
bilingüe y bidireccional. Además de la generación de texto a partir de RDF, se
incluyó la tarea inversa de extracción estructurada (\textit{text-to-RDF}). El
conjunto de datos incluyó inglés y ruso.

Las métricas empleadas en la tarea de generación RDF $\rightarrow$ texto fueron
BLEU, METEOR, chrF++, TER, BERTScore y BLEURT, este último únicamente para el
inglés~\cite{CastroFerreira2020}. El repositorio oficial del desafío se
encuentra disponible en: \url{https://github.com/WebNLG/challenge-2020}.

\subsection{WebNLG 2023 Shared Task}

La edición de 2023 tuvo como objetivo principal evaluar sistemas de generación
en lenguas con recursos limitados. Se incorporaron maltés, irlandés, bretón,
galés y ruso. A diferencia de las ediciones anteriores, esta se presentó en el
\textit{workshop} MM-NLG 2023 (Praga)~\cite{Cripwell2023}.

Las métricas empleadas fueron BLEU, METEOR, chrF++, TER y BERTScore. El
conjunto de datos se encuentra accesible en:
\url{https://huggingface.co/datasets/webnlg/challenge-2023}.

La Tabla~\ref{tab:webnlg-ediciones} sintetiza las tres ediciones por año,
idiomas, tareas, métricas y sede.

\begin{table}[ht]
	\centering
	\footnotesize
	\begin{tabular}{lp{1.6cm}p{2.7cm}p{3.0cm}p{2.7cm}}
		\hline
		\textbf{Edición} & \textbf{Año}                                                                           & \textbf{Idiomas}                      & \textbf{Tareas}
		                 & \textbf{Sede / Métricas}                                                                                                                                 \\
		\hline
		WebNLG           & 2017                                                                                   & inglés                                & RDF $\rightarrow$ texto
		                 & INLG, Santiago de Compostela; BLEU, METEOR, TER~\cite{Gardent2017}                                                                                       \\
		WebNLG+          & 2020                                                                                   & inglés, ruso
		                 & RDF $\rightarrow$ texto y \textit{semantic parsing} (texto $\rightarrow$ RDF)
		                 & INLG; BLEU, METEOR, chrF++, TER, BERTScore, BLEURT (sólo en)~\cite{CastroFerreira2020}                                                                   \\
		WebNLG           & 2023                                                                                   & bretón, galés, irlandés, maltés, ruso
		                 & RDF $\rightarrow$ texto en lenguas con recursos limitados
		                 & MM-NLG, Praga; BLEU, METEOR, chrF++, TER, BERTScore~\cite{Cripwell2023}                                                                                  \\
		\hline
	\end{tabular}
	\caption{Síntesis de las ediciones de WebNLG por idioma, tarea, sede y
		métricas de evaluación.}
	\label{tab:webnlg-ediciones}
\end{table}

\section{Estructura del dataset}

El conjunto de datos WebNLG está compuesto por instancias individuales
denominadas \textit{entries}. Cada \textit{entry} relaciona un conjunto de
tripletas RDF, generalmente extraídas de DBpedia, con una o varias
realizaciones en lenguaje natural. Esta organización permite evaluar sistemas
de los tipos \textit{data-to-text} y
\textit{text-to-data}~\cite{Gardent2017,CastroFerreira2020}. La descripción del
formato que sigue se apoya en la documentación oficial del consorcio WebNLG y
en el material descriptivo de la edición bilingüe de
2020~\cite{WebNLGGitHub,CastroFerreira2020}.

\subsection{La entry como unidad básica}

Cada \textit{entry} constituye un ejemplo independiente y posee cinco atributos
principales:

\begin{itemize}
	\item \textbf{category}: categoría semántica de DBpedia.
	\item \textbf{eid}: identificador del ejemplo, único dentro de una misma categoría y
	      partición (\textit{split}).
	\item \textbf{size}: número de tripletas RDF, entre 1 y 7.
	\item \textbf{shape}: representación estructural del grafo RDF.
	\item \textbf{shape\_type}: tipo de estructura del grafo.
\end{itemize}

El atributo \textit{size} permite medir la complejidad semántica de la oración,
mientras que \textit{eid} garantiza su identificación dentro de su grupo
estructural.

\subsection{Representación estructural: shape y shape\_type}

Los conjuntos de tripletas RDF se organizan como árboles. La estructura del
grafo se describe mediante el atributo \textit{shape}, que emplea una notación
anidada similar al formato Newick.

Adicionalmente, cada estructura se clasifica según el atributo
\textit{shape\_type}:

\begin{itemize}
	\item \textbf{Chain}: el objeto de un triple actúa como sujeto del siguiente.
	\item \textbf{Sibling}: varios triples comparten el mismo sujeto.
	\item \textbf{Mixed}: combinación de los patrones anteriores.
	\item \textbf{NA}: valor reservado para \textit{entries} que contienen una única
	      tripleta (\textit{size}\,=\,1), en las que la clasificación estructural no es
	      aplicable.
\end{itemize}

Esta categorización permite analizar el impacto de la complejidad estructural
en el rendimiento de los modelos.

\subsection{Contenido interno de una entry}

Cada \textit{entry} contiene tres bloques fundamentales:

\begin{itemize}
	\item \textbf{originaltripleset}: tripletas RDF extraídas directamente de DBpedia.
	\item \textbf{modifiedtripleset}: versión normalizada y corregida para entrenamiento.
	\item \textbf{lexs}: realizaciones en lenguaje natural (\textit{lexicalisations}).
\end{itemize}

Las lexicalizaciones poseen un identificador (\textit{lid}), un atributo de
idioma (\textit{lang}) y un comentario de calidad (\textit{comment}),
generalmente ``good'' o ``toFix''.

\subsection{Diferenciación entre triples originales y modificados}

La coexistencia de ambas versiones responde a una separación funcional. Los
triples originales preservan la trazabilidad con DBpedia, facilitando análisis
semánticos y alineaciones externas. Los triples modificados, en cambio,
garantizan uniformidad terminológica y estructural, reduciendo inconsistencias
léxicas.

Para tareas de entrenamiento y evaluación automática se recomienda emplear
exclusivamente la versión modificada.

\subsubsection{Ejemplo ilustrativo}
A continuación se presenta una \textit{entry} correspondiente al conjunto
bilingüe inglés--ruso publicado en WebNLG+ 2020~\cite{CastroFerreira2020}:

\begin{verbatim}
<entry category="Airport"
       eid="1"
       shape="(X (X))"
       shape_type="NA"
       size="1">

  <originaltripleset>
    <otriple>
      Aarhus_Airport | cityServed | "Aarhus, Denmark"@en
    </otriple>
  </originaltripleset>

  <modifiedtripleset>
    <mtriple>
      Aarhus_Airport | cityServed | "Aarhus, Denmark"
    </mtriple>
  </modifiedtripleset>

  <lex lang="en" lid="Id1">
    The Aarhus is the airport of Aarhus, Denmark.
  </lex>

  <lex lang="ru" lid="Id1">
    Орхус - аэропорт города Орхус, Дания.
  </lex>
</entry>
\end{verbatim}

\section{Extensión al español}

De manera análoga a la incorporación del ruso en la edición de 2020, es posible
extender el conjunto de datos al español mediante la inclusión de nuevas
lexicalizaciones alineadas:

\begin{verbatim}
<lex lang="es" lid="Id1">
  Aarhus es el aeropuerto de Aarhus, Dinamarca.
</lex>

<lex lang="es" lid="Id2">
  El aeropuerto de Aarhus da servicio a la ciudad de Aarhus, Dinamarca.
</lex>
\end{verbatim}

Asimismo, pueden añadirse enlaces auxiliares de alineación léxica:

\begin{verbatim}
<link direction="en2es">
  Aarhus Airport | includes | Aeropuerto de Aarhus
</link>
\end{verbatim}

Esta ampliación permite construir un subconjunto multilingüe coherente,
manteniendo la compatibilidad estructural con el \textit{benchmark} original.
